%%
%% This is a placeholder manuscript shell for WWW'27 (Research Track — Security and Privacy).
%% Refresh this manuscript before final submission by syncing with docs/paper/manuscript.tex
%% from the research repository.
%%
\documentclass[sigconf,anonymous,review]{acmart}

\acmConference[WWW '27]{The Web Conference 2027}{May 10--14, 2027}{Dublin, Ireland}
\acmYear{2027}
\acmDOI{}
\acmISBN{}

\title{Prompt Injection Through Web Content: The Span-Only Language Contrast Is Not Scored}

\author{Anonymous Author(s)}
\affiliation{%
  \institution{Anonymous Institution}
  \country{Anonymous Country}
}
\email{anonymous@example.org}

\begin{abstract}
The security question for autonomous Web agents is whether an untrusted instruction injected into retrieved third-party Web content causes an agent to issue unrequested tool calls, and evades detection when the user prompt remains in English but the retrieved language changes. A bounded evaluation across 1,160 InjecAgent episodes with locally served open-weight models shows attack success is near zero across languages (English 1.7\% and 0.0\%), falsifying pre-registered hypothesis H1 under a floor effect. No cross-lingual execution gap is claimed. Across 57,912 scored inferences, pre-registered hypothesis H2 is also falsified: open prompt injection detectors do not exhibit cross-lingual recall collapse at English-calibrated thresholds. Instead, multilingual encoders exhibit severe calibration asymmetry that elevates raw injection scores for non-English inputs, inflating false positives on benign text (up to 67.7\% false alarms in Chinese and 15.7\% in Arabic). Language-conditioned thresholds unmask evasion collapse, dropping Arabic recall from 54.0\% to 7.0\%. Subword tokenisation fertility reaches 4.33x in non-Latin scripts, identifying an architectural correlate for calibration drift. Defence-in-depth principles for secure Web agent architectures follow directly from these empirical findings.
\end{abstract}

\begin{CCSXML}
<ccs2012>
   <concept>
       <concept_id>10002978.10003022</concept_id>
       <concept_desc>Security and privacy~Software and application security</concept_desc>
       <concept_significance>500</concept_significance>
       </concept>
   <concept>
       <concept_id>10002951.10003260</concept_id>
       <concept_desc>Information systems~World Wide Web</concept_desc>
       <concept_significance>500</concept_significance>
       </concept>
 </ccs2012>
\end{CCSXML}

\ccsdesc[500]{Security and privacy~Software and application security}
\ccsdesc[500]{Information systems~World Wide Web}

\keywords{prompt injection, Web agents, security and privacy, retrieved content}

\begin{document}

\maketitle

\section{Introduction}
Autonomous Web agents interact with untrusted third-party Web content to satisfy user requests. In typical workflows, an agent retrieves external HTML, text, or API payloads and incorporates these spans into its reasoning context. When retrieved spans contain adversarial instructions, prompt injection occurs. This paper investigates whether changing the language of the untrusted retrieved span compromises safety mechanisms calibrated on English text.

\section{Threat Model and Web Agent Architecture}
We define the threat model for autonomous Web agents operating in multilingual environments. The user issues instructions in English, while the external retrieval mechanism ingests third-party spans in multiple languages (English, Chinese, Arabic, Spanish). The adversary aims to trigger unauthorised tool invocation or exfiltration without raising detector alarms.

\section{Experimental Design and Registered Hypotheses}
We formulate two pre-registered hypotheses:
\begin{itemize}
    \item \textbf{H1 (Agent Execution Gap):} Attack success rate increases when prompt injection spans are translated into non-English languages.
    \item \textbf{H2 (Detector Evasion Gap):} Open injection detectors calibrated on English show significant recall degradation on non-English injection payloads.
\end{itemize}
All evaluation parameters, datasets (InjecAgent), and metric definitions are pre-registered.

\section{Empirical Findings}
Evaluation over 1,160 InjecAgent episodes across locally served open-weight models shows attack success is near zero across all languages, falsifying H1 under an experimental floor effect. Across 57,912 detector inferences, open detectors exhibit calibration asymmetry rather than recall collapse: raw injection scores are elevated across all non-English text, yielding benign false alarms up to 67.7\% in Chinese and 15.7\% in Arabic. Applying language-conditioned thresholds causes Arabic recall to drop to 7.0\%. Subword tokenisation fertility analysis indicates up to 4.33x token expansion in non-Latin scripts.

\section{Web Agent Defence-in-Depth}
Based on these findings, we formulate defence-in-depth principles for Web agent architectures:
\begin{enumerate}
    \item Strict isolation between user control channels and retrieved third-party data spans.
    \item Multi-threshold calibration that accounts for language-dependent representation geometry.
    \item Tool-level authorization gates independent of input text classifiers.
\end{enumerate}

\section{Related Work}
We contextualise our results within recent research on prompt injection, indirect injection via Web retrieval, and multilingual robustness in transformer-based architectures.

\section{Conclusion}
The span-only language contrast reveals significant calibration asymmetry in multilingual injection detectors, rather than simple cross-lingual recall failure. These results mandate architectural isolation in Web agents rather than sole reliance on text classification guards.

%%
%% References and Appendices to be linked from research repository
%%
\bibliographystyle{ACM-Reference-Format}
\bibliography{references}

\end{document}
